\section{Workflow Implementation}

The workflow consists of \TcpAqmModule{}, \TcpAqmBenchmark{}, and Python runner and analysis scripts. The contrib module defines typed C++ objects for experiment and topology configuration. These objects load YAML or JSON files through \texttt{yaml-cpp}, validate the experiment and topology sections, map TCP and queue-disc names into enum values, store scalar values as ns-3 \texttt{Time}, \texttt{DataRate}, and \texttt{QueueSize} objects, and derive omitted bottleneck delays from the configured base RTT.

The C++ benchmark example resolves TCP aliases and TypeId names, configures the selected TCP variants per node, and keeps DCTCP-specific alpha tracing attached only when the resolved TypeId is DCTCP. This lets a campaign use short names such as \texttt{cubic} and \texttt{dctcp}, or full names such as \texttt{ns3::TcpBbr}, without changing the scenario source. The same binary can also be launched with \texttt{--configFile}, allowing a versioned configuration file to select a single-bottleneck or two-bottleneck topology and to control flow start times. For mixed-flow runs, the second flow can receive a bounded start-time jitter that is driven by the configured ns-3 random-run value.

The runner expands a named experiment matrix and invokes \texttt{./ns3 run} for each configuration. Each run receives an explicit \texttt{rngRun} value and stores a compact set of run artifacts:\RunArtifactPaths{}

\begin{itemize}
  \item Run metadata in \MetadataFile{}, including parameters, ns-3 commit, branch state, working-tree status, SHA-256 hashes of every benchmark and configuration source file that contributed to the run, a snapshot of ns-3-relevant environment variables (\texttt{NS\_LOG}, \texttt{NS\_GLOBAL\_VALUE}, \texttt{LD\_LIBRARY\_PATH}, and similar), return code, resolved trace files, and wall-clock time.
  \item \CommandFile{}, containing the exact ns-3 command.
  \item \StdoutFile{} and \StderrFile{}.
  \item Raw \TraceGlob{} trace files.
\end{itemize}

When the working tree contains modifications to the benchmark example, the contrib module sources, or any configuration file referenced by the run, the runner prints a warning. The per-file SHA-256 hashes still pin the exact source bytes used by that run, so a dirty tree does not silently invalidate the artifact.

The analyzer reads each run directory and writes a campaign-level \SummaryCsv{} and \SummaryAggregateCsv{}. The summaries include throughput, RTT, queueing delay, congestion window, drop count, mark count, and, for mixed-flow runs, second-flow throughput, total throughput, and Jain fairness (defined for the two-flow case used in this study). The analyzer clamps a configured warmup threshold below the first-flow start time so that pre-flow zero samples cannot deflate mean throughput. A dependency-free SVG plotter creates quick figures from the aggregate table for inspection before final plotting.

The contrib module includes a validation executable that loads a YAML or JSON file and prints the normalized configuration. This provides a lightweight check that configuration files are syntactically valid, semantically complete, and produce the intended derived topology values before a campaign is launched.

\begin{figure}[t]
  \centering
  \begin{tikzpicture}[
    node distance=0.6em and 0.6em,
    stage/.style={
      draw=black!55,
      line width=0.4pt,
      fill=black!4,
      rounded corners=3pt,
      align=center,
      text width=0.86\columnwidth,
      inner sep=5pt,
      font=\scriptsize,
    },
    header/.style={font=\scriptsize\bfseries},
    arrow/.style={
      -{Latex[length=2.2mm,width=1.8mm]},
      line width=0.7pt,
      draw=black!75,
    },
  ]
    \node[stage] (config) {%
      \textbf{Inputs}\\[1pt]
      YAML/JSON configs and named matrix\\
      \scriptsize TCP pair, queue, RTT, ECN, topology, \texttt{rngRun}};
    \node[stage, below=of config] (runner) {%
      \textbf{Python runner}\\[1pt]
      expands matrix, writes per-run \ConfigJson{}, invokes ns-3};
    \node[stage, below=of runner] (ns3) {%
      \textbf{\TcpAqmModule{} example}\\[1pt]
      \TcpAqmBenchmark{} loads typed config, builds topology, emits traces};
    \node[stage, below=of ns3] (runs) {%
      \textbf{Run directories}\\[1pt]
      \MetadataFile{}, command, stdout/stderr, raw trace files};
    \node[stage, below=of runs] (analysis) {%
      \textbf{Analysis}\\[1pt]
      \SummaryCsv{}, aggregate CSVs, derived evaluation tables, SVG figures};
    \draw[arrow] (config) -- (runner);
    \draw[arrow] (runner) -- (ns3);
    \draw[arrow] (ns3) -- (runs);
    \draw[arrow] (runs) -- (analysis);
  \end{tikzpicture}
  \caption{Reproducible TCP/AQM benchmark workflow. Each stage's output is the input of the stage below.}
  \Description{Vertical workflow diagram with five stages: inputs, Python runner, ns-3 benchmark example, run directories, and analysis; each stage is connected to the next by a downward arrow.}
  \label{fig:workflow}
\end{figure}
